\section{Resolution System}
\label{sec:resolution}

The \ms{} algorithm is parameterised by a choice of \emph{fine resolution}
and \emph{coarse resolution}.
The resolution system~\citep{dellaLibera2026resolution} defines three
options, all sharing the same \texttt{derive\_partition} helper via GEOID
prefix matching.

\subsection{GEOID Structure and Prefix Matching}

US Census Bureau GEOIDs encode geographic hierarchy in their digit structure.
For the census tract level, the 11-digit GEOID has the form:
\begin{center}
\texttt{SS CCC TTTTTT}
\end{center}
where \texttt{SS} (2 digits) is the state FIPS code, \texttt{CCC} (3 digits)
is the county FIPS code, and \texttt{TTTTTT} (6 digits) is the tract code.
Block group GEOIDs extend this with one additional digit: \texttt{SS CCC TTTTTT B}.

The \texttt{derive\_partition} helper takes a list of fine GEOIDs and a
prefix length, and groups fine units by their GEOID prefix:

\begin{definition}[\texttt{derive\_partition}]
\label{def:derive-partition}
Given a set of fine-unit GEOIDs $\{g_v : v \in V\}$ and a prefix length
$\ell \in \{5, 11\}$, the partition map $f_\ell : V \to C_\ell$ is defined
by
\[
  f_\ell(v) = g_v[0{:}\ell]
\]
where $g_v[0{:}\ell]$ denotes the first $\ell$ characters of the GEOID.
The coarse unit set $C_\ell = \{g_v[0{:}\ell] : v \in V\}$ is the image
of $f_\ell$.
\end{definition}

The coarse adjacency graph $G_{C_\ell}$ is then derived by
Definition~\ref{def:coarsening} using $f_\ell$ as the coarsening map.

\subsection{Three Resolution Options}

Table~\ref{tab:resolution-options} summarises the three options.

\begin{table}[h]
\centering
\caption{Multi-scale MCMC resolution options.}
\label{tab:resolution-options}
\begin{tabular}{lllll}
\toprule
Option & Fine level & Coarse level & GEOID prefix & Status \\
\midrule
A & Block group & Tract  & 11 digits & Defined \\
B & Tract       & County & 5 digits  & Implemented \\
C & Block group & County & 5 digits  & Defined \\
\bottomrule
\end{tabular}
\end{table}

\paragraph{Option A: Block group $\to$ Tract.}
Fine units are block groups (12-digit GEOID); coarse units are census tracts
(11-digit GEOID).
The prefix map uses length 11 to group block groups by their parent tract.
This option provides the finest spatial resolution available in the census
data hierarchy and is useful for states with very heterogeneous tract
populations.
Option A is defined but not yet implemented; it requires block-group
adjacency data (not currently loaded by the census pipeline).

\paragraph{Option B: Tract $\to$ County (implemented).}
Fine units are census tracts (11-digit GEOID); coarse units are counties
(5-digit GEOID = state FIPS + county FIPS).
The prefix map uses length 5 to group tracts by their parent county.
This is the implemented option and the focus of this paper.
Option B requires no extra data files: county adjacency is derived from the
existing tract graph as described in Section~\ref{sec:algorithm}.

\paragraph{Option C: Block group $\to$ County.}
Fine units are block groups; coarse units are counties.
The prefix map uses length 5 (same as Option B).
Option C skips the intermediate tract level and operates at the extremes of
the census hierarchy.
It is defined but requires block-group adjacency data.

\subsection{Shared Infrastructure}

All three options share the same \texttt{derive\_partition} code path.
The only difference is the input GEOID list (tract or block group) and the
prefix length (5 or 11).
The county adjacency derivation (Proposition~\ref{prop:county-adj-exact})
applies to Options B and C; for Option A, the tract adjacency is derived
analogously from the block-group graph.

The \texttt{rebalance} subroutine is identical across all three options:
it operates on the fine plan $\pi : V \to [k]$ and the fine adjacency graph
$G$, regardless of the coarse resolution.
The coarse tolerance $\epsilon_c = 3 \epsilon_f$ is also shared.

\subsection{Data Requirements}

\begin{table}[h]
\centering
\caption{Data requirements by option.}
\label{tab:data-reqs}
\begin{tabular}{llll}
\toprule
Option & Fine adjacency & Coarse adjacency & Extra files? \\
\midrule
A & Block group (not loaded) & Derived from fine & Yes \\
B & Tract (loaded) & Derived from tract & No \\
C & Block group (not loaded) & Derived from fine & Yes \\
\bottomrule
\end{tabular}
\end{table}

Option B is the only option that requires no extra data files beyond what
is already loaded by the standard redistricting pipeline.
The county graph is derived in $O(|E|)$ time at the start of each \ms{} run
and cached for the duration of the run.
For Texas ($|E| \approx 15{,}600$ tract edges), this derivation takes under
1 millisecond.
